\documentclass[12pt,a4paper]{article}

% Russian language support for XeLaTeX/LuaLaTeX
\usepackage{fontspec}
\usepackage{polyglossia}
\setdefaultlanguage{russian}
\setotherlanguage{english}

% Use fonts that support Cyrillic
\setmainfont{DejaVu Serif}
\setsansfont{DejaVu Sans}
\setmonofont{DejaVu Sans Mono}

\usepackage{geometry}
\geometry{margin=1in}

% Document metadata
\title{Тест русского языка}
\author{Docker Container}
\date{\today}

\begin{document}

\maketitle

\section{Введение}
Это тестовый документ для проверки поддержки русского языка в LaTeX Alpine Docker контейнере.

\section{Русский алфавит}
\subsection{Строчные буквы}
а б в г д е ё ж з и й к л м н о п р с т у ф х ц ч ш щ ъ ы ь э ю я

\subsection{Прописные буквы}
А Б В Г Д Е Ё Ж З И Й К Л М Н О П Р С Т У Ф Х Ц Ч Ш Щ Ъ Ы Ь Э Ю Я

\section{Примеры текста}
\begin{itemize}
    \item Привет, мир!
    \item Москва — столица России
    \item LaTeX поддерживает русский язык
    \item Математика: функция, интеграл, производная
    \item Программирование: алгоритм, структура данных
\end{itemize}

\section{Математические формулы}
В русском тексте можно использовать математические формулы:
\begin{equation}
\int_{0}^{\infty} e^{-x^2} dx = \frac{\sqrt{\pi}}{2}
\end{equation}

И inline формулы: $\sum_{i=1}^{n} i = \frac{n(n+1)}{2}$

\section{Смешанный текст}
\begin{english}
This is English text mixed with Russian.
\end{english}

А это снова русский текст после английского.

\section{Компиляция}
Для компиляции используйте:
\begin{verbatim}
xelatex test-russian.tex
\end{verbatim}
или:
\begin{verbatim}
lualatex test-russian.tex
\end{verbatim}

\textbf{Примечание:} Этот документ требует XeLaTeX или LuaLaTeX для корректного отображения кириллицы.

\end{document}